\documentclass[10pt]{article}
\usepackage[margin=20mm]{geometry}
\usepackage{amsmath,amssymb,mathtools,fontspec,unicode-math,xurl,hyperref}
\IfFontExistsTF{DejaVu Serif}{\setmainfont{DejaVu Serif}}{\setmainfont{Latin Modern Roman}}
\setmathfont{latinmodern-math.otf}
\setlength{\parindent}{0pt}\setlength{\parskip}{6pt}\sloppy
\title{Technical verification of A topological gap theorem for the pi-two systole of positive scalar curvature three-manifolds}
\author{AI-assisted technical review}
\date{}
\begin{document}
\maketitle
Abstract: The main inverse-mean-curvature-flow argument, the topological nontriviality of its spherical components, the exact gap constant, and the doubling argument are supported by the checks below. No Major mathematical defect is established. The background rigidity equivalence is stronger than the cited result and has an explicit counterexample; it should use a round universal cylinder. The neck counting and smoothing estimates require minor bookkeeping corrections, with complete choices given here. These changes preserve the main conclusions. Scientific adjudication remains pending.\par
\section*{Exact source and assessment}
The source is Kai Xu, A topological gap theorem for the \(\pi_2\)-systole of positive scalar curvature 3-manifolds, M25-0167, arXiv:2307.01922v3, 27 October 2024, 13 pages [source]. The assignment associates it with DOI 10.1215/00127094-2024-0061 (https://doi.org/10.1215/00127094-2024-0061). The supplied author version is the exact object reviewed; identity with the journal-final article or a later revision is not asserted. All page numbers below are one-based physical PDF pages and agree with the printed numbering.\par
The lead read the complete supplied text and actually inspected all 13 source pages. Three native helpers independently covered weak-flow preliminaries, the main topology and decomposition, and the doubling and gluing construction. This includes every stated theorem, claim, definition and proof, with explicit checks of covering invariance, orientability, initial-data conditions, weak time regularity, the terminal time, long-neck constants and the smoothing parameters. The external works were checked only in the specific scopes stated in Section 6, not represented as full independent reviews.\par
The supported assessment is Minor: the principal new theorems remain supported after bounded clarifications. The confirmed counterexample concerns the introductory rigidity formulation, not the topological gap. The paper itself leaves sharpness, equality and further gaps open; this review does not turn those open questions into defects.\par
\section*{The main topological and flow argument}
Let \(A_0=\operatorname{sys}_{\pi_2}(M,g)\) and \(\lambda=\min_M R_g\). If \(\lambda\le0\), the desired upper bound is immediate; only \(\lambda>0\) requires the positive-scalar-curvature classification. The oriented closed three-manifold is then a connected sum of spherical space forms and \(S^2\times S^1\). The finite fundamental-group case has \(\pi_2=0\). The infinite virtually cyclic cases are \(S^2\times S^1\) and \(\mathbb{RP}^3\#\mathbb{RP}^3\), both excluded by the topological covering hypothesis. The remaining free products have exponential growth. This is the growth input to the proper weak-flow existence theorem, not a geometric assumption on an arbitrary universal cover. More explicitly, the normal-cover consequence in [proper] gives isoperimetric lower bounds \(cv^{2/3}\) for small volume and \(cv^{1/2}\) for large volume. These make the reciprocal profile integrable at zero and the profile divergent at infinity, exactly the existence theorem’s hypotheses. Infinite deck translates of a small ball also verify the infinite-volume assumption.\par
\subsection*{Covering and the initial sphere}
Every sphere map lifts through a covering, and the induced map on based \(\pi_2\) is an isomorphism. Local isometries preserve the area of each lift, so systoles agree even though the corresponding sets of free homotopy classes need not be in bijection.\par
The instruction to remove all \(\mathbb{RP}^3\) factors by one double cover can be made explicit. Send the generator of every \(\mathbb Z/2\) free-product factor to the nonzero element of \(\mathbb Z/2\), and send all other factors and free generators to zero. The kernel has index two and intersects every conjugate of those \(\mathbb Z/2\) factors trivially. Kurosh decomposition and prime uniqueness give a cover with no \(\mathbb{RP}^3\) prime factor. An embedded projective plane in an oriented three-manifold has a regular neighborhood equal to a punctured \(\mathbb{RP}^3\); it would give such a factor. The Meeks–Yau minimizer can therefore be taken to be an embedded sphere, rather than a twofold projective plane. Its lift \(\widetilde\Sigma\) separates the simply connected universal cover.\par
Both sides of that sphere are noncompact. Otherwise one side \(W\) is compact and simply connected with boundary \(S^2\). Relative duality and the cohomology exact sequence give \(H_2(W;\mathbb Z)=H^1(W,\partial W;\mathbb Z)=0\); Hurewicz then makes its boundary null-homotopic in the universal cover, a contradiction.\par
\subsection*{Excision and outward minimization}
Start the proper weak flow with a one-sided collar in the unused side of \(\widetilde\Sigma\), restrict it to the other side, and cap off the initial sphere by a disk \(D\). The metric extends smoothly over the cap and the resulting manifold is complete. The variational definition is local outside the initial set, so the restriction, extended negatively in \(D\), is a proper solution with \(E_0=D\). No scalar-curvature lower bound is needed inside this arbitrary cap.\par
A bounded least-area envelope of the original collar gives an envelope of the capped disk. This follows by replacing its part on the retained side with any competing set and comparing perimeters; the collar fixes the trace at the separating sphere. If the envelope differs from \(D\), its boundary is stable minimal and represents \([\widetilde\Sigma]\) in the universal cover. Its components are two-sided and orientable. Stability with constant test function and \(\lambda>0\) gives \[0<\frac{\lambda}{2}|\Gamma|
 \le \int_\Gamma K=2\pi\chi(\Gamma),\] so each is spherical. Nonzero total homology forces at least one component to be nonzero in \(\pi_2\), with area at least \(A_0\). The original disk is therefore outward minimizing. This justifies the exact area law \[|\Sigma_t|=e^t A_0.
 \tag{1}\]\par
\subsection*{The homological control of every spherical component}
The flow’s connected sublevel sets and absence of compact exterior components give \[H_3(N\setminus E_0,\Sigma_t;\mathbb Z)=0.\] To see the two contributions separately, relative duality gives \(H^0(\overline E_t\setminus E_0,\Sigma_0;\mathbb Z)=0\), since every component meets the initial boundary, and \(H_c^0(N\setminus E_t;\mathbb Z)=0\), since no exterior component is compact. Excision and the exact sequence then make \(H_2(\Sigma_t;\mathbb Z)\to H_2(N\setminus E_0;\mathbb Z)\) injective. Noncompactness of the unused side similarly makes \[H_2(\widetilde M_2;\mathbb Z)\longrightarrow H_2(\widetilde M;\mathbb Z)\] injective. A spherical component is therefore nontrivial in \(\pi_2(\widetilde M)\), by Hurewicz. Although the weak flow gives \(C^{1,\alpha}\) surfaces, smooth isotopic approximation with convergent areas transfers the smooth immersion definition of systole to these spheres.\par
Each sphere consequently has area at least \(A_0\). Equation (1) shows there is at most one when \(t<\log2\). Every other orientable component has nonpositive Euler characteristic, hence \[\chi(\Sigma_t)\le2\qquad(0<t<\log2).
 \tag{2}\] This direct count also covers the case in which the first splitting time \(T\) defined on page 8 is infinite. The proof need not assume that its defining set is nonempty or that the infimum is attained.\par
\section*{The weak time calculation and the constant}
Put \(\mathcal H(t)=\int_{\Sigma_t}H^2\). Dropping nonnegative dissipation terms in the Geroch inequality gives \[\mathcal H(t_2)\le\mathcal H(t_1)+
 \int_{t_1}^{t_2}
 \left(4\pi\chi(\Sigma_s)-\lambda A_0e^s
                         -\tfrac12\mathcal H(s)\right)\,ds.
 \tag{3}\] The applicable weak-flow result [Retained latex: \textbackslash{}cite[Theorem 5.7]\{HI\}] allows the initial minimal surface and does not require a connected evolving surface, asymptotic flatness, or nonnegative scalar curvature inside the cap.\par
The paper correctly avoids differentiating \(\mathcal H\). For fixed \(t\), define the absolutely continuous function \(F(s)=\int_s^t\mathcal H(r)\,dr\). Rearranging (3) with \(t_1=s,t_2=t\) yields, almost everywhere, \[\frac{d}{ds}(e^{s/2}F(s))
 \le 4\pi e^{s/2}\int_s^t\chi(\Sigma_r)\,dr
 -\lambda A_0e^{s/2}(e^t-e^s)-e^{s/2}\mathcal H(t).\] Fubini and integration give the coefficient \((4/3)e^{3t/2}-2e^t+2/3\) in (2.5). Substitution into (3) proves \[\mathcal H(t)\le e^{-t/2}\mathcal H(0)
 +4\pi e^{-t/2}\int_0^t e^{s/2}\chi(\Sigma_s)\,ds
 -\frac23\lambda A_0(e^t-e^{-t/2}).\] Thus the factor \(2/3\), all signs and the exponential weights in Corollary 2.5 are correct.\par
The initial sphere is minimal, so \(\mathcal H(0)=0\). Using (2) and \(\mathcal H(t)\ge0\), for \(0<t\le\log2\) one obtains \[0\le16\pi(1-e^{-t/2})
       -\frac23\lambda A_0(e^t-e^{-t/2}).\] The single endpoint value of the Euler characteristic does not affect the integral. With \(x=e^{t/2}\) this simplifies to \[\lambda A_0\le\frac{24\pi}{x^2+x+1}.\] The right side decreases with \(t\); at \(t=\log2\) it is \[\frac{24\pi}{3+\sqrt2}
 =\frac{24\pi(3-\sqrt2)}7
 =24\pi\frac{2-\sqrt2}{4-\sqrt2},
 \tag{4}\] exactly the claimed approximately \(5.44\pi\). As \(t\downarrow0\), the same expression tends to \(8\pi\). This is a check of the stated bound and its optimization within the available time interval, not a proof of sharpness.\par
\section*{The local theorem and its bounded bookkeeping repairs}
For each stable minimal boundary sphere with metric \(h\), let \(\phi>0\) be the first stability eigenfunction. The Jacobi operator can be written \[-\Delta_h+K_h-\tfrac12(R+|A|^2).\] Its nonnegative lowest eigenvalue implies \(\Delta_h\phi\le(K_h-\lambda/2)\phi\). The translation-invariant neck metric \(g_i=h_i+\phi_i^2dt^2\) therefore satisfies \[R_{g_i}=2K_{h_i}-2\phi_i^{-1}\Delta_{h_i}\phi_i\ge\lambda.\] Its slices are totally geodesic. Gluing two copies of the punctured sphere by the \(k\) necks gives \(N=\#^{\,k-1}(S^2\times S^1)\), with \(k-1\ge2\). A neck slice is nonseparating because the other necks still connect the two copies; it is nontrivial in \(\pi_2(N)\). The least-area such slice gives \(\operatorname{sys}_{\pi_2}(N)\le A_0\). The free fundamental group has no projective-plane obstruction, so a systolic minimizer is an embedded sphere.\par
\subsection*{Uniform neck exclusion}
A minimal sphere entering a neck is either a horizontal slice or crosses every height. Otherwise its height range has an interior extremum and the maximum principle against the minimal slices forces coincidence. Equivalently the height satisfies the appropriate weighted harmonic equation on the surface. The local monotonicity estimate supplies a positive area \(c_0\) in a fixed-width slab, independently of neck length.\par
The displayed index range on page 10 includes possible end slabs outside the permitted interior range. A precise replacement is: choose an integer \[m>A_0/c_0+2,\qquad T=4m,\] and use the \(2m-1\) centers \(t=4j\), \(|j|\le m-1\). The slabs \((t-2,t+2)\) lie inside the neck, are disjoint up to endpoints, and each contributes at least \(c_0\). A crossing minimizer would have area \((2m-1)c_0>A_0\), contradicting the comparison slice. Alternatively one can first reduce \(c_0\) and enlarge the printed \(T\); its precise value is not part of the theorem. The reference to Claim 1 in this count should be to Claim 3.\par
A minimizer avoiding all necks lies in one punctured \(S^3\) piece. If its class were zero there, it would bound a ball and be trivial in the double. Its nonzero absolute homology class has an area-minimizing integral representative in that piece. The minimal boundary is a barrier; any interior component is stable, and hypothesis (ii) makes its homology class zero. Removing it lowers area. Thus the minimizing representative consists of boundary components with nonzero multiplicity and has area at least \(A_0\).\par
\subsection*{Smoothing with explicit parameter budgets}
The conformal perturbation \(g'=e^{2u}g\), with \(u=0\) and \(\partial_\nu u>0\) on the boundary, preserves the boundary tensor and makes its outward mean curvature positive. A collar parametrized with normal speed \(\phi_i\) makes the extended neck metric equal to \(g'\) as a full tensor at the seam; it is totally geodesic on the neck side. The hypotheses of the Brendle–Marques–Neves gluing theorem [BMN] are therefore met. The normal collar can be chosen using the original metric, since \(u=0\) makes \(g'=g\) at the seam; thus the smooth structure can be fixed independently of the small perturbation.\par
Its construction also permits arbitrary \(C^0\) closeness to \(g'\). Writing the boundary difference as \(\widetilde g-g'=\rho T\), the outer perturbation is \(O(\Lambda^{-1})T\), while the inner perturbation is bounded by \[\|T\|_\infty
 \bigl(e^{-\Lambda^2}+\Lambda e^{-2\Lambda^2}\bigr).\] Both tend to zero. This verifies the additional \(C^0\) clause used by this paper; it is not an unsupported strengthening.\par
The two displayed bounds of size \(\varepsilon^3\) on pages 10–11 do not literally sum to \(\varepsilon^3\). Choose each \(C^0\) perturbation at most \(\varepsilon^3/4\) instead, accounting for the uniformly equivalent tensor norms. Choose the conformal scalar loss and the gluing loss each at most \(\varepsilon/2\). These choices give a smooth metric equal to \(g_i\) on every neck, with \[\|g_N'-g_N\|_{C^0(g_N)}\le\varepsilon^3,\qquad
 R_{g_N'}\ge\lambda-\varepsilon.\] For a two-dimensional surface, this tensor bound gives area distortion at most \((1-\varepsilon^3)^{-1}<1+\varepsilon^2\) when \(0<\varepsilon<1/100\). Consequently \(\operatorname{sys}_{\pi_2}(N,g_N')\ge(1-\varepsilon)A_0\). Take also \(\varepsilon<\lambda\), apply Theorem 1.1, and let \(\varepsilon\downarrow0\). Theorem 1.2 follows with the same (4).\par
\section*{The introductory rigidity equivalence needs correction}
Problem 1 (Minor: equation (1.2), page 1).  Equality in the old \(8\pi\) bound does not imply a Riemannian covering by a compact round product \(S^2\times S^1\).\par
The primary result [BBN] instead identifies the universal cover with a round \(S^2\times\mathbb R\), up to scale.\par
Here is a counterexample to the stronger wording in this version. Let \(R_\theta\) be a sphere rotation with \(\theta/(2\pi)\notin\mathbb Q\), and take the quotient of the unit round \(S^2\times\mathbb R\) by \[\gamma(x,t)=(R_\theta x,t+L),\qquad L>0.\] This is a closed oriented manifold diffeomorphic to \(S^2\times S^1\) with scalar curvature two. The sphere area form descends to a closed form of comass one. Any nonzero \(\pi_2\) map lifts to the universal cover and its sphere projection has nonzero integer degree \(d\); thus its area is at least \(4\pi|d|\). A slice attains \(4\pi\), so \(\operatorname{sys}_{\pi_2}\,\min R=8\pi\).\par
A covering from a compact round product would have finite degree. On the universal cover, every isometry preserves the positive and zero Ricci eigendistributions, hence is a product isometry. It cannot conjugate a pure translation into \(\gamma^n(x,t)=(R_{n\theta}x,t+nL)\) for any \(n\ge1\). Irrationality makes every such rotation nontrivial. This disproves precisely the claimed compact-product covering.\par
The correction does not affect Theorem 1.1: this example is topologically \(S^2\times S^1\) and is excluded by its hypothesis. The invariant \(w(M)\) discussion on page 4 also needs the convention \(\operatorname{sys}_{\pi_2}=+\infty\) for \(\pi_2=0\) if it is to include spherical space forms; the original definition was restricted to nonzero \(\pi_2\).\par
\section*{Alternative decomposition and verification limits}
Remark 3.2 is consistent with the main result after clarifying its compressed steps. A finite-index free subgroup is obtained by mapping the finite free-product factors to their direct product; Kurosh and prime decomposition give a finite cover \(\#^k(S^2\times S^1)\). Small bumpy perturbations preserve positive scalar curvature. The stability area bound and compactness prevent an infinite disjoint family of essential stable spheres for a bumpy metric: an accumulating family would produce a Jacobi field, contrary to bumpiness. Systole is continuous under uniform metric comparison, so the perturbation can be removed at the end.\par
To exclude a punctured piece with a remaining \(S^2\times S^1\) factor, choose an absolute \(H_2\) class whose image in \(H_2(U,\partial U)\) is nonzero and minimize among closed cycles. Boundary components vanish in the relative group, so a minimizer must have an interior component. This supplies the contradiction with maximality without relying on an unspecified relative minimizer that might have free boundary. A piece with one boundary component would be a ball and make that cutting sphere inessential. If every piece had at most two, all would be cylinders and the connected gluing graph a cycle, yielding \(S^2\times S^1\). Therefore a piece with at least three boundary components exists. The word triviality in the last contradiction of Step 2 on page 12 should be nontriviality.\par
The specific primary checks were the Huisken–Ilmanen growth formula and its treatment of the initial minimal surface and arbitrary endpoints (printed pages 405–408), Xu’s properness theorem and normal-cover example in arXiv:2307.00725v2 (16 September 2024), the Brendle–Marques–Neves Theorem 5 and its cutoff construction, and the Bray–Brendle–Neves rigidity statement and minimizer proposition. The original Coulhon–Saloff-Coste PDF did not yield readable content in the available retrieval; the inspected normal-cover consequence in Xu’s primary article suffices for the existence application. These are bounded reference checks; the full analytic proofs of every cited theorem, the entire bibliography, and later developments were not independently reproduced. The covering, homological, weak-time and constant calculations, the explicit rigidity counterexample, and the local repairs above were independently worked out. No author code or simulation was needed, and unavailable code is not a mathematical defect. Native helpers are not human reviewers; their backend identities were not independently verified. Author adoption is unknown. Scientific decision remains pending and human reviews remain empty.\par
9 K. Xu, A topological gap theorem for the pi-two-systole of positive scalar curvature 3-manifolds, arXiv:2307.01922v3, 27 October 2024. https://arxiv.org/abs/2307.01922v3. H. Bray, S. Brendle and A. Neves, Rigidity of area-minimizing two-spheres in three-manifolds, Comm. Anal. Geom. 18 (2010), 821–830, Theorem 1.1 and Proposition 2.2. https://doi.org/10.4310/CAG.2010.v18.n4.a6. G. Huisken and T. Ilmanen, The inverse mean curvature flow and the Riemannian Penrose inequality, J. Differential Geom. 59 (2001), 353–437, Theorem 5.7 and the cited level-set properties. K. Xu, Isoperimetry and the properness of weak inverse mean curvature flow, Calc. Var. Partial Differential Equations 63, 216 (2024). https://doi.org/10.1007/s00526-024-02832-3. S. Brendle, F. C. Marques and A. Neves, Deformations of the hemisphere that increase scalar curvature, arXiv:1004.3088v3, 24 November 2010, Theorem 5 and its construction. https://arxiv.org/abs/1004.3088v3.\par
\end{document}
